\documentclass[11pt]{article}
\usepackage[margin=25mm]{geometry}
\usepackage{fontspec,unicode-math}
\setmainfont{Latin Modern Roman}
\setmathfont{Latin Modern Math}
\usepackage{amsmath,amssymb,mathtools}
\usepackage{xurl,hyperref}
\hypersetup{colorlinks=true,urlcolor=blue,linkcolor=blue}
\setcounter{secnumdepth}{0}
\setlength{\parindent}{0pt}
\setlength{\parskip}{5pt}
\setlength{\emergencystretch}{3em}
\newcommand{\sref}[2]{\hyperlink{source:#1:#2}{[S#2]}}
\newcommand{\cataloguescope}[1]{\paragraph{Source question.}#1}
\newcommand{\floor}[1]{\left\lfloor#1\right\rfloor}
\begin{document}
\section{An FPAS for Constrained Bicriteria Spanning Trees}\label{statement}

\subsection{Definitions and mathematical statement}

\paragraph{Shared notation.}Standard number systems are denoted by $\mathbb N,\mathbb Z,\mathbb Q,\mathbb R,\mathbb C$; index ranges, computational models and specialized notation follow the cited problem source and the definitions below.


A PTAS has approximation ratio $1+\varepsilon$ for minimization (or $1-\varepsilon$ for maximization), with polynomial runtime for each fixed $\varepsilon>0$. A fully polynomial scheme is polynomial also in $\varepsilon^{-1}$.

In a bicriteria minimization problem each feasible solution has two costs. A cost pair is Pareto-optimal if no feasible pair is no larger in both coordinates and strictly smaller in one.

A spanning tree of a connected finite graph is an acyclic connected subgraph containing every vertex. Its cost is the sum of its edge costs.
\paragraph{Statement recorded in the supplied source.}
Does the constrained bicriteria minimum spanning-tree problem admit a fully polynomial approximation scheme? \sref{statement}{2}
\paragraph{Scope qualification.}
The constrained objective, budget constraint, and permitted budget violation are those of the report; bicriteria budget violation is not silently permitted.
\paragraph{Source status.}
The supplied catalogue labels this record partially solved. Its source-status metadata was checked against the supplied snapshot on 2026-10-11. This import does not independently establish current open status or mathematical correctness. The catalogue literature assessment is dated 2026-08-21.
\subsection{Short English statement}
Decide whether constrained spanning-tree optimization with two costs permits an approximation scheme polynomial in both input size and inverse accuracy.
\subsection{Sources}
\begin{itemize}
\item[\hypertarget{source:statement:1}{S1}] Ulam AI and the UnsolvedMath Contributors, UnsolvedMath: A Curated Collection of Open Mathematics Problems, catalogue entry. Catalogue curation and metadata: CC BY 4.0. The supplied source question is retained; definitions and the English summary are editorial additions.
\url{https://huggingface.co/datasets/ulamai/UnsolvedMath}.
\item[\hypertarget{source:statement:2}{S2}] Approximation Algorithms for NP-Hard Problems (2005), Oberwolfach Reports, DOI 10.4171/owr/2004/28, pp. 1461--1540
\url{https://doi.org/10.4171/owr/2004/28}.
\item[\hypertarget{source:statement:3}{S3}] fully polynomial bicriteria approximation scheme
\url{https://doi.org/10.1016/j.orl.2003.06.003}.
\item[\hypertarget{source:statement:4}{S4}] PTAS for constrained minimum spanning tree
\url{https://doi.org/10.1137/S0097539703426775}.
\textit{Underlying Oberwolfach Reports material retains its source-specific licence, including CC BY-SA 4.0 where applicable; see THIRD\_PARTY\_NOTICES.md.}
\end{itemize}
\paragraph{Preparation and review.}
Original question attributed to the cited source and UnsolvedMath contributors. Codex (GPT-6; exact serving revision not exposed) assisted with definitions, English summaries, source inspection and formatting. No new human mathematical authorship or independent proof verification is asserted. Catalogue identity checked against the complete supplied registry and existing repository statements; known flagged repetitions and obvious identical conjectures excluded. Full mathematical-equivalence/literature review has not been performed. Statement preparation and final review followed contribution prompts C and D; automated checks verify structure and identity, not open status.
\label{end:statement}
\end{document}
